% ---- pg-463 （书页 451） ----

类. 换而言之, 错误分类的机会或概率应该很小. 我们将会看到, 一个“最优”分类法则还必须具备另外一些特性.

可能有这样的情况: 某一个类别或总体比另一个类别有更大的出现概率, 因为两个总体中一个比另一个大得多. 例如, 通常财务上健康的企业要大大多于濒临破产的企业. 另一个例子是: 某一类鸢尾比另一类更常见. 一个最优法则则应该考虑这些“先验出现概率”. 如果我们确信财务困难且最终将破产的企业的出现(先验)概率非常小, 就应将一个随机选取的企业分类为非破产的, 除非数据压倒性地倾向于破产.

分类的另一个问题是代价. 假如将一个 $\pi_1$ 对象误分为 $\pi_2$ 同将一个 $\pi_2$ 对象分入 $\pi_1$ 相比是一个更严重的错误, 那么我们就应对前者更加谨慎. 作为一个例子, 未能诊断出某种致命的疾病比错误地诊断有病(而实际上没有)是更严重的错误. 在任何可能的情况下, 最优分类方法应考虑与错误分类相联系的代价.

设 $f_1(\mathbf{x})$ 与 $f_2(\mathbf{x})$ 为 $p\times1$ 向量随机变量 $\mathbf{X}$ 分别在总体 $\pi_1$ 和 $\pi_2$ 中的概率密度函数. 设 $\Omega$ 为样本空间——所有可能的观测值 $\mathbf{x}$ 的集合. 一个具有观测值 $\mathbf{x}$ 的对象必须被指定属于 $\pi_1$ 或 $\pi_2$. 令 $R_1$ 为 $\mathbf{x}$ 值的一个集合, 当 $\mathbf{x}$ 属于 $R_1$ 时, 我们把对象分类为 $\pi_1$; $R_2=\Omega-R_1$ 为剩余的 $\mathbf{x}$ 值的集合, 对这些 $\mathbf{x}$ 值, 我们把对象分类为 $\pi_2$. 由于任一对象必定属于且只属于 $\pi_1$ 和 $\pi_2$ 之一, 因此 $R_1$ 和 $R_2$ 互斥且穷尽所有可能性. 对 $p=2$, 我们可能有形如\hyperref[fig:11.2]{图 11.2} 所示的情况.

\begin{mvfigure}
\mvinc{pg-463}{0.126,0.358,0.468,0.481}
\phantomsection\label{fig:11.2}\mvcap{图 11.2\quad 两个总体的分类区域}
\end{mvfigure}

将 $\pi_1$ 中的对象错分到 $\pi_2$ 的条件概率 $P(2\mid1)$ 为:
\phantomsection\label{eq:11-1}
\begin{equation}
P(2\mid1)=P(\mathbf{X}\in R_2\mid\pi_1)=\int_{R_2=\Omega-R_1}f_1(\mathbf{x})d\mathbf{x}
\tag{11-1}
\end{equation}
类似地, 将 $\pi_2$ 中的对象错分到 $\pi_1$ 的条件概率 $P(1\mid2)$ 为:
\phantomsection\label{eq:11-2}
\begin{equation}
P(1\mid2)=P(\mathbf{X}\in R_1\mid\pi_2)=\int_{R_1}f_2(\mathbf{x})d\mathbf{x}
\tag{11-2}
\end{equation}
在\hyperref[eq:11-1]{式(11-1)}中, 积分符号表示由密度函数 $f_1(\mathbf{x})$ 在区域 $R_2$ 中所形成的体积. 类似的, \hyperref[eq:11-2]{式(11-2)}中的积分符号表示由 $f_2(\mathbf{x})$ 在区域 $R_1$ 中所形成的体积. \hyperref[fig:11.3]{图 11.3} 在单变量($p=1$)场合对此作了说明.

\begin{mvfigure}
\mvinc{pg-463}{0.298,0.702,0.707,0.806}
\phantomsection\label{fig:11.3}\mvcap{图 11.3\quad 对假想分类区域的错分概率($p=1$)}
\end{mvfigure}

令 $p_1$ 为 $\pi_1$ 的先验概率, $p_2$ 为 $\pi_2$ 的先验概率, $p_1+p_2=1$. 此时正确或错误地将对象分类的无条件概率可表示为先验概率与条件概率之积:
